\section{総括 — 3月、競争軸はモデル単体から実行条件へ広がった}
\label{sec:retrospective-synthesis}
\noindent\textbf{本号を横断すると、2026年3月はfrontier modelの更新が続く一方で、serving、security、multimodal、agent runtime、evaluationがそれぞれ独立した設計課題として厚みを増した月だった。重要なのは、これらを一つの「Agent化」へ回収するのではなく、AIをどう動かし、つなぎ、守り、測るかという実行条件の広がりとして読むことである。}
\medskip
\begin{multicols}{2}
\par\medskip\Needspace{5\baselineskip}
\subsection{Frontier ModelsとServing — 更新される能力をどの実行面へ載せるか}
GPT-5.4やGemini APIなどの更新では、model identityだけでなくtool use、computer use、multimodal、API surfaceが同時に動いた。一方、TensorRT-LLMやvLLM、routing、cache、communicationを扱うEvidenceは、modelをproduction workloadとして動かす層が独立した設計対象になっていたことを示す。両者は同じ技術ではないが、3月には「何ができるか」と「どう実行するか」を別々に追う必要が明確になった。 \cite{src-eea6230f612a,src-a24fd97eb3d0,src-7e670bf82538,src-d6cda62f02bb}

\par\medskip\Needspace{5\baselineskip}
\subsection{RuntimeとSecurity — action surfaceが増えるほど境界設計が要る}
Responses APIのcomputer environmentはshell、container、files、tools、stateを持つruntimeを示した。同じ月のsecurity側では、instruction hierarchy、prompt injection対策、tool privilege、chain-of-thought monitoringが並ぶ。ここから読めるのは、model outputをactionへ変える経路が増えるほど、入力のtrust、実行権限、監視を別々のcontrolとして設計する必要が強まるという構造である。 \cite{src-01572b4c34e9,src-18b7ee4fe5e4,src-3dec299964e6,src-5becf449a527}

\par\medskip\Needspace{5\baselineskip}
\subsection{Multimodal — Agentに吸収されない独立した設計軸}
LongCat-Nextはmodalitiesをshared discrete spaceへ寄せる方向を、Unify-Agentは外部evidenceでimage synthesisをgroundする方向を、DreamLiteはon-device generation/editingを統合する方向を示した。searchやagentic pipelineとの接点はあるが、3月のmultimodalはrepresentation、grounding、deploymentという独自の設計問題を持っており、Agentの一部としてだけ整理するには広すぎる。 \cite{src-ac36472a8f09,src-197775677c0d,src-b8369fd2c988}

\par\medskip\Needspace{5\baselineskip}
\subsection{Long-horizonとEvaluation — 一回の成功率だけでは測れない}
Beyond pass@1、Marco DeepResearch、MemFactoryは、長時間taskでreliability、verification、memoryがmodel外部の設計変数になることを示した。Paper WatchではMiroEvalやELT-Bench-Verifiedがprocess/outcomeやbenchmark qualityを問い直している。長く動くsystemほど、最終回答だけでなく途中過程と測定器そのものを評価対象にする必要がある、という問題意識が共通している。 \cite{src-0ec537ec1cc1,src-2b8b39e5cae3,src-42318ebdee75,src-5d84874f2d03,src-a27679e5527b}

\end{multicols}
\Needspace{0.34\textheight}
\sectionkicker{Cross-chapter synthesis}
\subsection*{各章はどうつながるか}
\addcontentsline{toc}{subsection}{各章はどうつながるか}
\begin{center}
\small
\renewcommand{\arraystretch}{1.16}
\begin{tabularx}{\linewidth}{@{}p{0.28\linewidth}X@{}}
\toprule
接続する章 & 本号で見える構造的な関係 \\
\midrule
Frontier Models $\leftrightarrow$ Inference \& Serving & model／APIの更新と、cache・routing・RPC・communicationを担う実行基盤が別層の競争として同時に厚くなった。 \cite{src-eea6230f612a,src-7e670bf82538,src-d6cda62f02bb} \\
Agents \& Runtime $\rightarrow$ Safety \& Security & computer environmentやtool useがaction surfaceを増やす一方、instruction hierarchy、privilege、monitoringがcontrol surfaceとして具体化した。 \cite{src-01572b4c34e9,src-18b7ee4fe5e4,src-5becf449a527} \\
Multimodal $\leftrightarrow$ Systems & multimodal側ではrepresentation／grounding／deploymentが進み、systems側ではdistributed servingやmemory processingが別の実行課題として現れた。 \cite{src-ac36472a8f09,src-361b6111931d,src-b6e517afaa35} \\
Long-horizon $\rightarrow$ Evaluation & task duration、memory、verificationが増えるほど、pass@1だけでなくprocess、outcome、benchmark qualityまで含む評価設計が必要になる。 \cite{src-0ec537ec1cc1,src-5d84874f2d03,src-a27679e5527b} \\
\bottomrule
\end{tabularx}
\renewcommand{\arraystretch}{1.0}
\end{center}
\normalsize
\begin{claimboundary}[この総括の境界]
この総括は本号で選択した一次資料・原論文とaccepted articleを横断した編集上のsynthesisである。vendor／project／authorが報告した性能、attack result、benchmark結果を独立再現済みの事実へ昇格させず、異なる評価条件の数値を直接比較しない。arXivで後日revisionされた資料は3月の初回公開時系列と区別し、後日の記述を3月時点へ遡及させない。
\end{claimboundary}
\subsection*{3月をどう位置づけるか}
\addcontentsline{toc}{subsection}{3月をどう位置づけるか}
2026年3月は「Agentの月」と一言でまとめるより、frontier modelの更新を続けながら、その外側のserving、security、multimodal、runtime、evaluationが互いに独立した技術領域として前景化した月と位置付ける方が、当月のEvidenceに忠実である。model capabilityだけでAI systemを説明することが難しくなり、「どう動かすか」「何につなぐか」「どこで止めるか」「どう測るか」が同時に重要になったことが、3月全体を貫く変化だった。 \cite{src-eea6230f612a,src-7e670bf82538,src-18b7ee4fe5e4,src-ac36472a8f09,src-01572b4c34e9,src-0ec537ec1cc1}
